% Transcription — fonds Alexandre Grothendieck, shelfmark 100, batch 2 (pages 21-40).
% Established from the facsimile archives/batches/100/batch-02.pdf (a mirror of
% https://grothendieck.umontpellier.fr/100.pdf), per the skill
% transcribe-grothendieck. The batch file carries no cover sheet: PDF page N
% is page 20+N of the folder (the Montpellier footer prints « 21 » ... « 40 »),
% and the archivists' pencil numbers bottom-left agree wherever legible
% (« 22 », « 25 », « 27 », « 29 », « 31 », « 33 », « 35 » ... « 40 »; on the
% two leaves scanned upside down, 21 and 23, the pencil number is top right
% of the scan and reads « 21 », « 23 » once turned).
%
% Pass: Opus 5.5 (claude-opus-5-5), 2026-09-23 — first pass, unchecked against the pages by a human.
%
% The folder sits in the inventory group "69-89" « Géométrie et topologie
% combinatoire », dated 1976-[vers 1986], whose subject does not match this
% folder's; \dating{} carries the folder's own date verbatim. Nothing is
% concluded here from the group.
%
% The batch alternates: his manuscript on the even pages 22-40 (blue ink),
% typed leaves and letters on the odd pages 21-39.
% Eleven of the twenty pages are transcribed: 21, 22, 24, 26, 28, 30, 32,
% 34, 36, 38, 40. Absent:
%   - pages 23, 25, 27, 31, 35, 37: leaves of the same typed English text on
%     quotients of schemes by group actions that batch 1 met on pages 5, 7,
%     9, 13, 15 (here « -0.16 bis- » with Lemma 0.3, « -0.4- » with
%     Definition 0.4, « -0.2- » with Definition 0.1, and preface leaves
%     « Preface », « ii », « iv », the last signed « Harvard University,
%     Oct. 27, 1963 »), not by him and unnamed on the leaves. None carries a
%     mark that is certainly his: 23 and 27 have only stray typing ticks.
%     23 was scanned upside down;
%   - page 33, preface leaf « iii » of the same typescript, with three small
%     English corrections in blue ink (« that it is » circled into the line,
%     « developed » over a struck « considered », an inserted « r ») in a
%     rounded hand that does not look like his. Left out on that doubt; if a
%     reader judges the hand his, it belongs in the apparatus of a revision;
%   - page 29, a typed letter to him signed L. Illusie, « Paris, le 24
%     Février 1968 », about an exposé and a CNRS application, struck by two
%     pencil strokes; page 39, a typed letter « Cher Dieudonné » dated
%     « 7.4.1967 » on EGA errata (12.1.1, 12.1.3, 15.5.6.1, 15.6.3-15.6.4),
%     struck by diagonal strokes. Neither bears on these notes. They are
%     named only so that the gap is explained; nothing is inferred from their
%     dates about the date of the manuscript.
% Page 21 is kept: it is a leaf « -0.9- » of the same typescript (scanned
% upside down; turned here), and it carries his pencil on the typed text and,
% at its foot, a few lines of his notes in ink on the subject of the folder.
% Following batch 1 (pages 5, 7, 9, 15), the typed text is summarised in a
% French \note, not recomposed, and his marks are given as \marginal.
%
% What the batch is. A single run, unnumbered by him, headed on page 22 by
% « Lemme 1 »: descent of abelian schemes along a proper epimorphism, and
% the algebraisation of formal abelian schemes that it comes down to.
%   21     His notes: A — A', T_p(A''_1) ≃ T_p(A''_2), S — S' ⇉ S'',
%          S − T = T ⊔ T, T ⊔ T ⊔ T; pencil « loc. » twice before « finite type », and
%          « présentation » in the margin.
%   22     Lemme 1: f: S' → S an epimorphism of (loc. noetherian) schemes, X
%          formally unramified over S; then Hom_S(S,X) → Hom_S(S',X) ⇉
%          Hom_S(S'',X) is exact. Proof: g' factors through S'_1 = Spec(𝒜),
%          𝒜 = Ker(f_*O_S' ⇉ h_*O_S''), and S'_1 → S is finite radiciel
%          (« EGA III 3 »). Corollaire 1: bijectivity of Hom_S(S,X) →
%          Hom_S(S',X) when f is integral radiciel, or X → S loc. of f. type.
%   24     A reduction to the local case, much of it framed and struck;
%          Corollaire 2: a proper epimorphism of loc. noetherian schemes is
%          a descent morphism for relative abelian schemes (Hom_gr(A,B) is
%          representable and unramified, apply Lemma 1). Corollaire 3: the
%          same for p-divisible groups, reduced to showing Hom_gr(A,B)
%          formally unramified.
%   26     Effectivity (f an effective epimorphism): a) reduction to S local;
%          b) to S local complete, via fpqc descent « de Raynaud » and a
%          four-row array of 1-descent / 0-descent; c) a struck case
%          dim S = 0; d) announced: a formal abelian scheme over Spf R.
%   28, 30 d) Construction of the formal abelian scheme via the formal
%          moduli M of A_0 over T = Spf W(k), maps φ': Ŝ' → M, exactness of
%          Ŝ ← Ŝ' ⇇ Ŝ''; reduction to the Theorem (page 30): for R complete
%          local noetherian and  a formal abelian scheme over Spf R, (i) Â
%          algebraisable ⇔ (ii) algebraic after a proper surjective S' → S
%          ⇔ (iii) algebraic over the normalisation of every one-dimensional
%          integral quotient B of R. (i)⇒(ii), (ii)⇒(iii) proved.
%   32, 34 (iii)⇒(ii'), with ii' = a finite surjective S' → S: the schemes
%          N_n = NS_{A_n/S_n}, P_n of polarisations glued into P̂ = ∐ P̂^α,
%          Z_α = Im(P^α → S); then reductions (S reduced, étale extension,
%          trivial residue extensions); a Lemma (algebraicity checked on the
%          irreducible components, by gluing) and a Lemma (checked on S_red).
%   36, 38 Proof of the last Lemma by the infinitesimal theory (obstruction
%          in H², indeterminacy in H¹ of 𝒯 ⊗ J, « the same » formally and
%          algebraically); a Proposition: in a local noetherian S, countably
%          many closed parts whose union contains every x with
%          dim closure(x) ≤ 1 contain a finite subfamily covering S; its
%          Corollary (rare Z_i miss some x with dim closure = 1), via
%          V[[T_1..T_n]] and V-points in m^n.
%   40     Remarque: the same for a proper flat formal scheme X̂ with
%          H⁰(X_0, O) ≃ k: NS_{X_n/S_n}, P_n = NS^+, and four equivalent
%          conditions for X̂ to be algebraisable and projective; a long
%          marginal NB, and a struck last line.
% Page 22 names the lemma « Lemme 1 »; pages 24-26 cite « lemme 1 » and
% « cor. 2 »; page 30 cites « le th. de comparaison »; page 34 cites
% « le lemme suivant »; page 36 opens « Il reste à prouver la Proposition »
% before stating it. No reference to a page number.
%
% Hand. Fast; statements and formulas carry, connectives often do not and
% are left \ill{}. Pages 24 and 26 carry large framed blocks struck by wavy
% or diagonal lines; given as one \struck{} with one \note each (hand.md §6).
%
% Notation fixed for the batch. S, S', S'' = S' ×_S S', S''' as he writes;
% h: S'' → S; O_S as \mathcal{O}_S (his underlined O); 𝒜 (his script A) as
% \mathcal{A}; formal completions with \hat (Ŝ, Â', T̂); « Hom_gr » with
% \underline where he underlines; NS, Pic underlined as \underline{\mathrm{NS}},
% \underline{\mathrm{Pic}}; his underlined words in prose as \emph. Page 26's
% array of descent rows and page 28's sketch are set as arrays, their
% strokes being unheaded or overlapping; page 22's square is a tikz-cd.
%
% Cross-references out of the batch: none explicit. Page 26's « de
% Raynaud » (fpqc descent of abelian schemes) is the reference batch 1
% page 6 also makes (« d'après la réduction de Raynaud »), and page 21's
% notes use batch 1's T_p(A); no link is asserted beyond the shared words.
% Page 40 ends on a struck line; batch 3 (41-52) is not read here.
\documentclass[11pt,a4paper]{article}
% Le préambule commun à toutes les transcriptions du fonds Grothendieck.
%
% Il tient en une idée : **l'incertitude se compose, elle ne se résout pas.**
% Une transcription qui lisse un mot douteux a détruit la seule chose pour
% laquelle on la faisait — un lecteur doit pouvoir savoir, à chaque endroit,
% ce qui était sur le feuillet et ce qui a été ajouté. Les six macros qui
% suivent sont donc l'appareil critique, et non de la décoration.
%
% Les mêmes commandes sont reconnues par `scripts/render.mjs`, qui produit la
% vue de lecture du site. Une macro ajoutée ici doit l'être aussi là-bas,
% faute de quoi le rendu échouera bruyamment — ce qui est voulu : un rendu
% silencieusement faux est pire qu'un rendu absent.

\ProvidesPackage{grothendieck}[2026/08/08 Transcriptions du fonds Grothendieck]

\usepackage{amsmath,amssymb,amsthm}
\usepackage{mathrsfs}
\usepackage{stmaryrd}
% Les diagrammes commutatifs sont partout dans ce fonds : c'est la langue
% même de la géométrie algébrique de Grothendieck, et une transcription qui
% les rendrait en prose ne transcrirait rien.
\usepackage{tikz-cd}
% Français par défaut — c'est la langue des feuillets — mais l'anglais est
% chargé aussi : la traduction et le résumé sont des documents anglais, et la
% typographie française (espaces avant les deux-points) y serait une faute.
% Ces documents basculent par \selectlanguage{english} en tête de corps.
\usepackage[english,french]{babel}
\usepackage{fontspec}
\usepackage{geometry}
\geometry{a4paper,margin=25mm,marginparwidth=28mm}
\usepackage{marginnote}
\usepackage{xcolor}
% ulem, not soul. `soul`'s \st and \ul are built on a character-by-character
% scan that XeLaTeX's font machinery breaks: the document fails with an
% undefined control sequence rather than degrading. ulem does the same two jobs
% and is Unicode-engine safe. `normalem` stops it hijacking \emph, which the
% transcriptions use for Grothendieck's own emphasis.
\usepackage[normalem]{ulem}
\usepackage{hyperref}
\hypersetup{colorlinks=true,linkcolor=black,urlcolor=black,pdfborder={0 0 0}}

% --- Métadonnées du lot -----------------------------------------------------
% Reprises telles quelles par le script de rendu, qui les affiche en tête de
% la vue de lecture. Elles fixent ce que le fichier prétend couvrir : sans
% elles, une transcription flotte sans point d'attache dans un fonds de seize
% mille feuillets.

\newcommand{\folder}[1]{\gdef\grothFolder{#1}}
\newcommand{\batch}[1]{\gdef\grothBatch{#1}}
\newcommand{\leaves}[2]{\gdef\grothFirst{#1}\gdef\grothLast{#2}}
\newcommand{\foldertitle}[1]{\gdef\grothFoldertitle{#1}}
\gdef\grothFolder{?}\gdef\grothBatch{?}\gdef\grothFirst{?}\gdef\grothLast{?}\gdef\grothFoldertitle{}

% --- Le résumé --------------------------------------------------------------
%
% Il ouvre la lecture modernisée, sous le titre « Résumé », et il est composé
% en petit. Ce n'est pas un ornement : c'est le seul passage du document qui
% ne soit pas des mathématiques, et un lecteur doit pouvoir le reconnaître
% avant de l'avoir lu — puis le sauter s'il connaît déjà le sujet, ou s'y
% arrêter s'il ne le connaît pas. Le filet qui le suit marque la reprise du
% corps.

\newenvironment{resume}
  {\small\setlength{\parskip}{0.45em}}
  {\par\medskip\hrule\medskip}

% --- Mots-clés ---------------------------------------------------------------
%
% En anglais, en clôture du résumé de la lecture modernisée : le vocabulaire
% moderne sous lequel chercher ce que les feuillets construisent. Le manifeste
% du site (`npm run manifest`) les extrait pour étiqueter la cote entière —
% cette ligne est la source unique des tags, il n'y a pas de fichier à côté.
\newcommand{\keywords}[1]{\par\smallskip\noindent
  {\footnotesize\textsc{Keywords} --- \textit{#1}}\par}

% --- L'appareil critique ----------------------------------------------------

% Le numéro de feuillet, en marge. C'est l'ancre qui permet de régler un
% désaccord sur une formule en regardant *un* feuillet plutôt qu'une chemise
% de six cents. Le site s'en sert aussi pour tourner les pages du fac-similé
% quand on fait défiler la transcription.
\newcommand{\leaf}[1]{%
  \marginnote{\footnotesize\color{gray!70}\textsf{#1}}%
  \label{leaf:#1}\ignorespaces}

% Illisible : on ne devine pas. Un mot inventé qui se lit comme les autres est
% le pire résultat possible.
\newcommand{\ill}{\textcolor{red!60!black}{[\,\ldots\,]}}

% Lecture douteuse : le mot est donné, et signalé comme incertain.
\newcommand{\uncertain}[1]{\uline{#1}}

% Ajout de l'éditeur — une lettre manquante, un mot sous-entendu.
\newcommand{\add}[1]{\textcolor{blue!50!black}{[#1]}}

% Biffé par Grothendieck lui-même. Ce qu'il a rayé fait partie du document :
% une rature dit souvent où la pensée a changé de direction.
\newcommand{\struck}[1]{\sout{#1}}

% Note du transcripteur, sur l'état matériel du feuillet.
\newcommand{\note}[1]{\par\smallskip{\small\color{gray!60!black}\textit{[#1]}}\par\smallskip}

% Note marginale de Grothendieck, distincte du corps du texte.
\newcommand{\marginal}[1]{\par\smallskip\noindent\hspace{1em}%
  {\small\color{gray!60!black}#1}\par\smallskip}

% --- Titre ------------------------------------------------------------------

\AtBeginDocument{%
  \begin{center}
    {\large\bfseries Fonds Alexandre Grothendieck --- cote n\textsuperscript{o}~\grothFolder}\\[2pt]
    {\itshape\grothFoldertitle}\\[4pt]
    {\small feuillets \grothFirst--\grothLast\ (lot \grothBatch)}\\[2pt]
    {\footnotesize\color{gray!60!black}Transcription établie d'après le fac-similé numérisé
     par l'Université de Montpellier.\\
     Les lectures incertaines sont soulignées, les illisibles notées [\,\ldots\,],
     les ajouts entre crochets.}
  \end{center}
  \vspace{1em}}

\endinput


\folder{100}
\batch{2}
\pages{21}{40}
\dating{[à partir de 1967-1968]}
\watermark{Édition de démonstration}
\foldertitle{Bonne réduction des variétés abéliennes via bonne réduction de Barsotti-Tate Tp (A) (base quelconque) : notes manuscrites (s.d.), copies de tapuscrit annoté (s.d.).}

\begin{document}

\section{[Notes au pied d'un feuillet dactylographié]}
\note{titre de l'éditeur, entre crochets.}

\page{21}
\note{Le feuillet est une page dactylographiée en anglais, paginée
« -0.9- », du même texte sur les quotients d'un schéma par l'action d'un
groupe que le lot 1 rencontre aux pages 5 à 15 ; elle n'est pas de lui et
n'est pas recomposée ici. Elle porte la fin d'une discussion des quotients
catégoriques et géométriques (le cas d'un schéma normal algébrique sur un
corps), le point (3) (quotient catégorique universel, condition (iv) de la
définition 0.6, surjectivité, critère pour un quotient géométrique
universel) et le point (4) (projection $G \to S$ universellement ouverte :
plate et de type fini, le cas d'un groupe algébrique, cf.\ SGA 4, le cas
d'un $S$ irréductible normal noethérien, « Chevalley »). Ses marques au
crayon sont données ci-dessous ; au pied du feuillet, tête-bêche par rapport
au texte dactylographié, quelques lignes de sa main à l'encre, traversées de
trois longs traits courbes.}

\note{au point (4), « loc. » est ajouté au crayon, deux fois, au-dessus de
« of finite type » ; le mot « type » est entouré, et en marge, en face :}
\marginal{loc.} \quad \marginal{loc.} \quad
\marginal{\uncertain{présentation}}

\[
  A - A' \qquad \struck{\ill{}}\ A' \quad \struck{A''_1}\,,\ A''_2 \qquad
  T_p(A''_1) \simeq T_p(A''_2)
\]
\[
  S - S' \rightrightarrows S'' \qquad
  S - T = T \sqcup_{\ill{}} T = T \sqcup T \qquad
  T \sqcup T \sqcup T
\]
\[
  \uncertain{\mathfrak{Y}} - Y' \qquad Y'' = Y' \sqcup Y'
\]
\note{dans le dernier « $T \sqcup T \sqcup T$ », les deux derniers $T$ sont
soulignés ; l'indice du premier $\sqcup$ ne se lit pas.}

\section{[Descente des schémas abéliens par un épimorphisme propre]}
\note{titre de l'éditeur, entre crochets ; le feuillet 22 commence par
« Lemme 1 ».}

\page{22}
\textbf{Lemme 1.} Soit $f : S' \to S$ un épimorphisme \uncertain{propre} de
schémas \uncertain{[loc.\ noethériens]}, et \struck{$X$} $X$ un schéma
\struck{non ramifié} \emph{formellement} non ramifié sur $S$. Alors la suite
\[
  \mathrm{Hom}_S(S, X) \longrightarrow \mathrm{Hom}_S(S', X)
  \rightrightarrows \mathrm{Hom}_S(S'', X)
\]
est exacte, où $S'' = S' \times_S S'$.
\note{« propre » et « loc. » sont écrits au-dessus de la ligne ; les
crochets autour de « loc. noethériens » sont de lui, et renvoient à la note
entourée de la marge.}

\marginal{$[\;]$ inutile si $X$ loc.\ de t.f.\ ?}

\textbf{N.B.} Le fait que $f$ est un épimorphisme s'exprime par la
condition que l'homomorphisme
\[
  \mathcal{O}_S \longrightarrow f_*(\mathcal{O}_{S'})
\]
est \emph{injectif}.

Soit
\[
  g' : S' \longrightarrow X
\]
tel que $g' \mathrm{pr}_1 = g' \mathrm{pr}_2$, il faut montrer que $g'$ se
factorise par $S$. Tout d'abord il se factorise par $S$ en tant
qu'application continue, soit $g_0$, et
$f^{-1}(g_0^{-1}(\mathcal{O}_X)) \simeq g'^{-1}(\mathcal{O}_X) \to
\mathcal{O}_{S'}$ définit un hom
\[
  g_0^{-1}(\mathcal{O}_X) \longrightarrow f_*(\mathcal{O}_{S'}) ;
\]
l'hyp.\ $g' p_1 = g' p_2$ s'exprime par le fait que
$g_0^{-1}(\mathcal{O}_X)$ s'envoie dans
\[
  \mathcal{A} = \mathrm{Ker}\bigl(f_*(\mathcal{O}_{S'}) \rightrightarrows
  \struck{\ill{}}\, h_*(\mathcal{O}_{S''})\bigr)
\]

\begin{tikzcd}
  & & X \arrow[d] \\
  S'' \arrow[r, "p_1"] \arrow[r, "p_2"'] \arrow[rr, bend right=40, "h"'] & S' \arrow[r, "f"] \arrow[ur, "g'"] & S \arrow[u, dashed, bend right=40, "g_0"']
\end{tikzcd}

\note{croquis dans la marge gauche, en regard des lignes précédentes ; les
deux flèches $p_1$, $p_2$ sont parallèles ; la flèche $g_0$ est pointillée.}

donc $g'$ se factorise par
\[
  S' \longrightarrow \mathrm{Spec}(\mathcal{A}) = S'_1 \longrightarrow X .
\]
Or EGA III 3 indique aussitôt que $S'_1 \to S'$ est \uncertain{fini}
\emph{radiciel}, c'est d'ailleurs un épimorphisme car
$\mathcal{O}_S \hookrightarrow \struck{\ill{}}\ \mathcal{A} \subset
f_*(\mathcal{O}_{S'})$. On est donc ramené \ill{} un
\uncertain{épimorphisme} \uncertain{fini}, \ill{} à prouver le
\marginal{\emph{lemme}}
\note{« fini » est ajouté au-dessus de « radiciel » ; le « $S'$ » de
« $S'_1 \to S'$ » est peut-être un « $S$ ».}

\textbf{Corollaire 1.} Soit $f : S' \to S$ un épimorphisme
\struck{loc.\ de t.f.}\ \uncertain{fermé} de schémas \struck{\ill{}}, $X$
\struck{non ramifié} sur $S$ ; \struck{Alors si $f$ fini} si on suppose $f$
\struck{fini} \uncertain{entier} radiciel et \uncertain{$S$ loc.\ noeth.},
ou $X \to S$ loc.\ de t.f. Alors
\[
  \mathrm{Hom}_S(S, X) \longrightarrow \mathrm{Hom}_S(S', X)
\]
est \emph{bijectif}.
\note{la ligne est très reprise : « entier » en interligne au-dessus d'un
« fini » biffé, « et $S$ loc. noeth. » au-dessous ; une boucle renvoie
« si on suppose $f$ » à la fin de la ligne précédente. La lecture de
l'ordre exact des conditions est incertaine.}

\page{24}
Pour le voir, \struck{\ill{}} on est ramené \ill{} : \ill{} que \ill{}
\ill{} \uncertain{injectif}, \ill{} \uncertain{surjectif}. Soit
$g' : S' \to X$ \ill{} ; \uncertain{comme} $S' \rightrightarrows S$
\uncertain{est un conoyau}, prouvons que $g$ se factorise par $S$.
\struck{Or \ill{} $g$ \ill{}} \ill{} \uncertain{épimorphisme}, on est
ramené, \uncertain{localisant} au voisinage de $s$ \ill{} $S$, de $s'$
\uncertain{local}, \ill{}

\struck{\ill{} \ill{} $X' = X \times_S S'$ \ill{} $S'$, \ill{} \ill{} \ill{}
\ill{} $X'/S'$ est \uncertain{formellement net} \ill{} $X'$ ; \ill{}
$X' \to X$ \ill{} \ill{} \ill{} ; $g(S') = U$ \ill{} \ill{} \ill{} ;
\ill{} \ill{} \ill{} $x = g(s')$ \ill{} $\mathcal{O}_{X,x} \to
\mathcal{O}_{S',s'}$ \ill{} \uncertain{se factorise} par
$\mathcal{O}_{S,s}$ \ill{} ; \uncertain{Remplaçant} $X$ \ill{} \ill{}
voisinage \uncertain{affine} de $x$, on peut supposer $X$ affine, \ill{}
$X \to S$ \ill{} \ill{} $g'(S')$ \ill{} \ill{} \ill{} \ill{}}
\note{tout ce bloc est encadré et biffé d'un long trait ondulé et de traits
horizontaux ligne à ligne ; on n'en donne que ce qui se lit.}

\note{croquis dans la marge gauche :
$\mathcal{O}_{X,x}$ au-dessus de $\mathcal{O}_{S,s}$, flèche
$\mathcal{O}_{S,s} \to \mathcal{O}_{X,x}$ vers le haut,
$\mathcal{O}_{S,s} \to \mathcal{O}_{S',s'}$ vers la gauche, et une flèche
pointillée $\mathcal{O}_{X,x} \dashrightarrow \mathcal{O}_{S',s'}$ en
diagonale}

\uncertain{Remplaçant} $X$ par l'image \ill{} \ill{} \ill{} ; on peut
supposer \ill{} \ill{} $\mathcal{O}_{S,s}$ \ill{} \ill{}
\ill{} $\mathcal{O}_{S'}$ \ill{} \ill{} \ill{} $\mathcal{O}_{S'}$.
Alors $\mathcal{O}_{X,x}$ est donc \uncertain{entier} sur
$\mathcal{O}_{S,s}$, \ill{} $S' \to S$ \ill{}
\note{ces deux lignes, en partie reprises dans le bloc biffé, sont écrites
serrées sous lui ; leur liaison avec ce qui précède est incertaine.}

\textbf{Corollaire 2.} Soit $f$ un épimorphisme propre de schémas
localement noethériens. Alors $f$ est un morphisme de descente pour la
catégorie des schémas abéliens relatifs.
\note{un trait vertical en marge gauche encadre l'énoncé.}

Si $A$ et $B$ sont deux schémas abéliens sur $S$, on sait en effet que
$H = \underline{\mathrm{Hom}}_{\mathrm{gr}}(A, B)$ est représentable, et
est net sur $S$, donc il suffit d'appliquer le lemme 1.

\textbf{Corollaire 3.} Soit $p$ un nombre premier, et soit $f : S' \to S$
comme dans le cor.\ précédent. Alors $f$ est un morphisme de descente pour la
catégorie des \struck{gr} schémas en groupes $p$-divisibles.

L'argument \struck{pr} du lemme 1 nous ramène à prouver l'analogue du
corollaire 1, prenant $X = \underline{\mathrm{Hom}}_{\mathrm{gr}}(A, B)$,
pour $A, B$ deux groupes $p$-divisibles. Il faut prouver que $X$ est
formellement net.

\page{26}
Soit $f$ comme dans Cor.\ 2, (mais on suppose $f$ un épim.\
\uncertain{effectif}). On veut prouver que $f$ est un morphisme de descente
effective. Soit donc $A'$ schéma abélien sur $S'$, avec donnée de descente
relativement à $S' \to S$, il faut prouver que cette donnée est effective.

a) Réduction au cas $S$ local ($= \mathrm{Spec}\,R$) : triviale, par passage
à la limite.

\struck{\ill{} \ill{} \ill{}}

b) Réduction au cas $S$ local complet : \uncertain{critère} \ill{} de
descente fpqc de schémas abéliens \uncertain{de Raynaud}. Mais il faut
savoir que $S'_2 \to S_2$ est \uncertain{morphisme} de descente pour la
catégorie des schémas abéliens relatifs, et on ne peut directement appliquer
le \struck{lemme 1, ou} cor.\ 2, car $S_2$ n'est plus loc.\ noeth.\,! Mais on
s'en \uncertain{tire} \ill{} \ill{} \ill{} que \uncertain{puisqu'il est
vrai} \struck{\ill{}} \ill{} que (comme on utilise ce qui a été dit dans le
lemme 1) \struck{\ill{} est radiciel sur \ill{}}
\[
  \mathrm{Ker}\bigl(f_{2*}(\mathcal{O}_{S'_2}) \rightrightarrows
  h_{2*}(\mathcal{O}_{S''_2})\bigr)
\]
est \uncertain{radiciel} sur $\mathcal{O}_{S_2}$, \ill{} \ill{} \ill{}
\uncertain{relation} \ill{}
$\mathcal{O}_S \to \mathrm{Ker}\bigl(f_*(\mathcal{O}_{S'})
\rightrightarrows h_*(\mathcal{O}_{S''})\bigr)$ \ill{}

\[
  \begin{array}{ccccccccl}
    S & - & S' & = & S'' & \equiv & S''' & & \text{1-descente} \\
    \mid & & \mid & & \mid & & \mid & & \\
    S_1 & - & S'_1 & = & S''_1 & \equiv & S'''_1 & & \text{1-descente} \\
    \parallel & & \parallel & & \parallel & & \parallel & & \\
    S_2 & \underset{f_2}{-} & S'_2 & = & S''_2 & \equiv & S'''_2 & &
      \text{1-descente} \\
    \vert\vert\vert & & \vert\vert\vert & & \vert\vert\vert & &
      \vert\vert\vert & & \\
    S_3 & - & S'_3 & = & S''_3 & \equiv & S'''_3 & & \text{0-descente}
  \end{array}
\]
\note{tableau dans la marge gauche, en regard de b) ; ses traits n'ont pas
de pointes. Entre les deux dernières lignes, trois traits verticaux, le
signe « $\equiv$ » tourné ; le « $f_2$ » est écrit sous le premier trait de
la troisième ligne. Les indices de la dernière ligne sont peu lisibles.}

\marginal{\ill{} \uncertain{donnée} \ill{} \uncertain{grâce au} lemme 1,
\uncertain{donnée} \ill{} \ill{} \uncertain{en 0-descente}}

\struck{c) Cas $\dim S = 0$ \ill{} \ill{} $\mathrm{long}_A\, A'/A = 1$
\ill{} \ill{} \ill{}}
\note{le paragraphe c) est encadré et biffé.}

d) Construction d'un schéma \uncertain{(abélien)} formel sur
$\mathrm{Spf}\,R$, qui descend le schéma formel défini par $A'$
(\uncertain{complété}) \struck{\ill{} $A'_3$ \ill{}}. \struck{Immédiat par
c)} \uncertain{Facile}.

\page{28}
d) \emph{Construction d'un schéma abélien formel.}

\[
  \begin{array}{ccccccccc}
    & & & & \hat{A}' & & \hat{A}'' & & \hat{A}''' \\
    M & \longleftarrow & S_0 & \longleftarrow & S'_0 & \leftleftarrows &
      S''_0 & & \\
    \downarrow & & & & & & & & \\
    T & \longleftarrow & \hat{S} & \longleftarrow & \hat{S}' &
      \leftleftarrows & \hat{S}'' & \Lleftarrow & \hat{S}''' \\[1ex]
    & & S_0 & - & S'_0 & = & S''_0 & \equiv & S'''_0
  \end{array}
\]
\note{croquis dans la marge gauche ; les $\hat{A}$ sont écrits au-dessus des
colonnes, sans flèches.}

Notons d'abord que si $S_0 = \mathrm{Spec}\,k$ ($k$ corps résiduel) dans
$S'_0/S_0$ \ill{} une \uncertain{section}, donc $S'_0 \to S_0$ est un
\uncertain{morphisme} de descente effective universelle, et on trouve un
$A_0$ sur $S_0$ qui descend $A'_0$.

Commençons par prendre les complétés formels de $S'$, $S''$, $S'''$, $A'$,
$A''$, $A'''$ relativement aux fibres en $s$, on trouve \struck{\ill{}}
schéma abélien formel $\hat{A}'$ sur $\hat{S}'$, avec donnée de descente
par $\hat{S}' \to \hat{S}$. On se ramène à descendre dans la catégorie des
schémas abéliens formels. On introduit $T = \mathrm{Spf}\,W(k)$ ($W$ vecteurs
de Witt) et $M =$ schéma formel des modules de $A_0$ ($M$ construit comme
schéma formel). Un lemme facile \ill{} qui est pratiquement trivial si $S'$
fini sur $S$ \uncertain{[donc $\hat{S}'$, $\hat{S}''$, \ldots\ des
\ill{} noethériens]} dit que la donnée d'une extension formelle $\hat{A}'$
de $A'_0 = A_0 \otimes S'_0$ [resp.\ $\hat{A}''$ de $A''_0$] équivaut à la
donnée d'un $T$-morphisme $\varphi' : \hat{S}' \to M$ [resp.\
$\varphi'' : \hat{S}'' \to M$, \struck{$\hat{S}''' \to M$}]. La donnée de
descente \struck{\ill{}} sur $\hat{A}'$ \uncertain{telle que} l'on a sur
$A'_0$ équivaut à la condition $\varphi' \hat{p}_1 = \varphi' \hat{p}_2$.
Or on \ill{} aussitôt (\uncertain{utilisant} \ill{}, triviale encore si $S'$
fini sur $S$, et \uncertain{prenant} \ill{} th.\ de comparaison
formel-algébrique en général) que la suite
\[
  \hat{S} \longleftarrow \hat{S}' \leftleftarrows \hat{S}''
\]
\struck{\ill{}} \uncertain{dans} les schémas formels est exacte, donc il
existe $\varphi : \hat{S} \to M$ tel que $\varphi \hat{f} = \varphi'$. Cela
signifie qu'il y a \ill{} descente \ill{}

\page{30}
$\hat{A}'$ en un schéma formel abélien $\hat{A}$. Le th.\ de comparaison
alg.-formel nous ramène à prouver que $\hat{A}$ est algébrique. On est ainsi
ramené à prouver le

\textbf{Th.} Soit $R$ anneau local noeth.\ complet, $S = \mathrm{Spec}\,R$,
$\hat{A}$ un \uncertain{schéma abélien formel} sur $\hat{S} =
\mathrm{Spf}\,R$. Conditions équivalentes :

(i) $\hat{A}$ algébrisable, i.e.\ $\exists$ schéma abélien $A$ sur $S$,
tel que $\hat{A}$ soit isom.\ à son complété formel \struck{le long de
$A_s$}.

(ii) $\exists$ morphisme propre surjectif $f : S' \to S$ \struck{\ill{}}
tel que l'image inverse $\hat{A}'$ de $\hat{A}$ sur $\hat{S}'$ soit
algébrique.

(iii) Pour tout \struck{algèbre finie $B$ sur $A$} quotient intègre $B$ de
$A$, de dim $1$, si $B'$ est la normalisée de $B$, et \struck{\ill{}}
$\hat{T}' = \mathrm{Spf}\,B'$, $\hat{A}_{\hat{T}'}$ est algébrique.
\note{« quotient de $A$ » : il faut sans doute lire $R$, l'anneau de base ;
on garde la lettre du manuscrit. « schéma abélien formel » est une
correction en interligne.}

\marginal{\textbf{N.B.} \ill{} \ill{}}

\struck{N.B. \ill{} \ill{}}

\emph{Dém.} (i) $\Rightarrow$ (ii) trivial. (ii) $\Rightarrow$ (iii) On
introduit $T' = S' \times_S T$, on voit que $\hat{T}' \to \hat{T}$ « rend
$\hat{A}_{\hat{T}}$ algébrique ». Or on peut trouver, comme $T$ est un
trait complet \struck{et fini}, un autre trait $T_1$ fini sur $T$, et un
$T$-morphisme $T_1 \to T'$. Donc $\hat{A}_{\hat{T}_1}$ est algébrique, i.e.\
définit un schéma abélien $A_{T_1}$ sur $T_1$. Ce dernier est muni d'une
donnée de descente rel.\ à $T_1 \to T$ [grâce au th.\ de comparaison], il
suffit de voir que la donnée est effective --- ce qui est évident, car
$T_1 \to T$ est fid.\ plat fini, et $A_1$ est projectif.
\note{ici $T = \mathrm{Spec}\,B'$, sans que la page le redise.}

\page{32}
\struck{(iii) $\Rightarrow$ (i) (cas $R$ réduit)}

(iii) $\Rightarrow$ (ii') $\Bigl\{$ $\exists$ morphisme fini surjectif
$S' \to S$, tel que $\hat{A}_{\hat{S}'}$ soit algébrique.

Soit \struck{$N_n \neq P_n$} $N_n = \underline{\mathrm{NS}}_{A_n/S_n}$,
$P_n =$ \struck{\ill{}} \uncertain{l'ouvert induit} défini par les
polarisations \struck{\ill{}}

Les $P_n$ se recollent et forment un schéma formel $\hat{P}$ sur
$\mathrm{Spf}\,R = \hat{S}$, \struck{\ill{}} (réunion d'une famille
dénombrable de schémas finis nets)
\[
  P^{\alpha} \text{ sur } S ; \qquad
  \hat{P} = \coprod_{\alpha} \widehat{P^{\alpha}}
\]
Soit \struck{$Z$} $Z_{\alpha} = \mathrm{Im}(P^{\alpha} \to S)$.

L'hyp.\ (iii) signifie que pour tout $B$ comme \uncertain{ci-dessus}, il
existe un $S$-morphisme $T' \to P^{\alpha}$ \struck{$\mathrm{Spec}\,B'$}
pour un $\alpha$ \uncertain{convenable}, donc le pt générique $x$ de
$T = \mathrm{Spec}\,B$ est dans $Z_{\alpha}$. Donc
$Z = \bigcup_{\alpha} Z_{\alpha}$ contient tous les $x \in S$ tels que
$\dim \overline{\{x\}} = 1$. Par un lemme général, il s'ensuit qu'il y a une
suite finie des $Z_{\alpha}$ qui recouvrent $S$. Nous prenons pour $S'$ la
somme des $P^{\alpha}$ correspondants.

(ii') $\Rightarrow$ \struck{(i') (cas $R$ réduit)} \ill{}
$\hat{A}_{\hat{S}_1}$ est algébrique, si $S_1 = S_{\mathrm{red}}$. On peut
supposer ainsi $S$ réduit. Ensuite, faisons une extension étale
\uncertain{finie} $S_1 \to S$ (et \uncertain{redescendre}, grâce \ill{} ce
qui est possible par Raynaud) \ill{} \ill{} on peut supposer que

\page{34}
les extensions résiduelles de $S' \to S$ sont triviales
\struck{, et que les composantes irréductibles de $S$ sont géom.\
unibranches}. Alors pour tt composant connexe $S'_i$ de $S'$,
$S'_i \to S$ est une immersion \struck{\ill{}} fermée. Si $S$ est
irréductible, une des immersions fermées est un isom.\ (car la somme des
\ill{} est surjective) et on conclut. \struck{\ill{} on utilise ensuite un
\ill{}} \struck{Pour tt composant}

\textbf{Lemme.} Soit $\hat{X}$ schéma formel propre sur $\hat{S}$
($S = \mathrm{Spec}\,R$, $R$ anneau loc.\ noeth.\ complet), soit $S$
\ill{} réduit. Supposons que pour tt composant irréductible (muni
\uncertain{structure induite réduite}) $S_i$ de $S$,
$\hat{X} \times_{\hat{S}} \hat{S}_i$ soit algébrique. Alors $\hat{X}$ est
algébrique.

Facile par récurrence sur nb de composantes : $S = S' \cup S''$,
\uncertain{irréductibles}, on est ramené au cas où $\hat{X}_{\hat{S}'}$ et
$\hat{X}_{\hat{S}''}$ \struck{\ill{}} algébriques, $S'$ et $S''$ sous-schémas
fermés, et proviennent de $X'$ et $X''$. Alors $X'_T \simeq X''_T$
($T = S' \cap S''$) par th.\ de comparaison, et on recolle $X'$ et $X''$
par cet isom. On voit facilement que $X$ recollé marche.

\uncertain{(ii) $\Rightarrow$ (i)} : Cela résulte du lemme suivant
(\ill{} \ill{} résultat démontré \ill{} \ill{} en particulier \ill{} \ill{}
ab.\ relatif \ill{}) :
\note{l'étiquette « (ii) $\Rightarrow$ (i) » est lue telle quelle ; d'après
la page 32, il s'agit de passer de (ii') au cas réduit.}

\textbf{Lemme.} Soit $\hat{X}$ schéma formel propre \struck{lisse} sur
$\hat{S}$, où $S = \mathrm{Spec}\,R$, $R$ anneau local noeth.\ complet. Soit
$S' = S_{\mathrm{red}}$, et supposons que $\hat{X}_{\hat{S}'}$ soit
algébrique. Alors $\hat{X}$ l'est aussi.

\page{36}
On est ramené à prouver que si \struck{\ill{}} \ill{} \uncertain{un idéal
nilpotent}, $S_0 = V(J)$, $J^2 = 0$, et si $\hat{X}_{\hat{S}_0}$ est
algébrique, alors $\hat{X}$ l'est aussi. On utilise la théorie de
prolongement infinitésimal d'une structure \uncertain{lisse} $X_0/S_0$
\uncertain{vers} $\hat{X}/\hat{S}$, et

1°) L'obstruction à prolonger $X_0$ (resp.\ $\hat{X}_0$) en un $X$ (resp.\
$\hat{X}$) lisse sur $S$ (resp.\ $\hat{S}$) est dans « le même » groupe
\[
  H^2(X_0, \mathcal{T}_{X_0/S_0} \otimes_{\mathcal{O}_{S_0}} J) \simeq
  H^2(\hat{X}_0, \mathcal{T}_{\hat{X}_0/\hat{S}_0}
  \otimes_{\mathcal{O}_{\hat{S}_0}} \hat{J})
\]
et moyennant cette identification, l'obstruction est la même.

2°) L'indétermination pour \uncertain{prolonger} est dans « le même »
groupe
\[
  H^1(X_0, \mathcal{T} \otimes J) \simeq H^1(\hat{X}_0, \hat{\mathcal{T}}
  \otimes \hat{J}) .
\]
\note{le lemme de la page 34 est pour un $\hat{X}$ propre, « lisse » y étant
biffé ; l'argument de cette page suppose la lissité.}

Il reste à prouver la

\textbf{Proposition.} Soit $S$ un \uncertain{schéma local}
\uncertain{(noethérien)} \struck{anneau local noeth.\ complet}
\struck{\ill{}}, $(Z_i)_{i \in I}$ une famille \uncertain{dénombrable}
\struck{suite} de parties fermées \struck{\ill{}} de $S$, telles que
$Z = \bigcup Z_i$ contienne tt $x \in S$ tel que
$\dim \overline{\{x\}} \leqslant 1$. Alors il existe une partie finie $J$ de
$I$ telle que $\bigcup_{i \in J} Z_i = S$.
\note{« schéma local » et « (noethérien) » sont écrits au-dessus d'une
première rédaction biffée ; la fin de cette ligne biffée, « de dim $\geqslant
1$ », l'est aussi. Le $I$ de « $i \in I$ » surcharge une autre lettre.}

Soit $S = \mathrm{Spec}\,R$. Si $\dim S \leqslant 1$, c'est trivial (\ill{}
\ill{} \uncertain{$S$ est fini}, $Z = S$). Si $n = \dim S \geqslant 2$,
\ill{} \ill{} l'hyp.\ \ill{} que \ill{} \ill{} \ill{} \ill{}
\uncertain{disons} qu'il existe un anneau \ill{} \uncertain{de val.\
discrète} complet $V$, \struck{\ill{}} tel que $R$ soit fini sur
\[
  V[[T_1, \ldots, T_n]] .
\]
\uncertain{Notons} que \ill{} \ill{} \ill{} \ill{}
\note{la fin du feuillet, sous une grande accolade barrée d'un trait, ne se
lit pas ; au-dessus, « \uncertain{annulés} ».}

\page{38}
\textbf{Corollaire.} Supposons que les $Z_i$ soient \emph{rares} dans $S$.
Alors il existe un $x \in S$, avec $\dim \overline{\{x\}} = 1$, tel que
\[
  x \notin \bigcup Z_i .
\]
\note{un « le » isolé en tête du feuillet, au-dessus de « Corollaire ».}

\uncertain{Remplaçant} $S$ par $T = \mathrm{Spec}\,V[[T_1, \ldots, T_n]]$,
\struck{\ill{}} on voit que les $T_i = \mathrm{Im}(Z_i \to T)$ sont
\uncertain{rares}, et il suffit de trouver un $x \in T$, avec
$\dim \overline{\{x\}} = 1$, tel que $x \notin \bigcup T_i$. (En effet, il
existe alors \ill{} un $y \in S_x$ au-dessus de $x$ \ill{}
\uncertain{nécessairement} $\dim \overline{\{y\}} = 1$.) On
\uncertain{remarque} qu'il existe \struck{\ill{}} un $x \in T$ provenant
d'un $V$-hom
\[
  V[[T_1, \ldots, T_n]] \longrightarrow V . \qquad \struck{\ill{}}
\]
En effet, l'ens.\ de ces hom.\ s'identifie à $\mathfrak{m}^n$
($\mathfrak{m} = \mathrm{rad}(V)$). Il a une structure d'espace
\uncertain{métrique} \struck{\ill{}} complet. Montrons que pour
\struck{\ill{}} partie fermée rare $Z$ de $T$, l'ens.\ des $\varphi$ tels que
$V(\mathrm{Ker}\,\varphi) \subset Z$ est rare. En effet, on peut supposer
\struck{$Z$ est défini} $Z = V(f)$,
$f \in V[[T_1, \ldots, T_n]]$, $f \neq 0$, et
$|V(\mathrm{Ker}\,\varphi)| \subset |V(f)|$ (\uncertain{car}
$V(\mathrm{Ker}\,\varphi)$ est intègre \ill{} \uncertain{réduit})
équivaut \uncertain{(comme} \ill{}
\[
  V(\mathrm{Ker}\,\varphi) \subset V(f) \ \ldots \qquad
  f \in \mathrm{Ker}\,\varphi \ \ldots \qquad \varphi(f) = 0 ;
\]
i.e.
\[
  f(\varphi_1, \ldots, \varphi_n) = 0 .
\]
Il est immédiat que c'est une partie fermée rare.
\note{« rares » est souligné dans l'énoncé. La parenthèse à droite de
« $Z = V(f)$ » est ajoutée d'une encre plus pâle, comme la ligne des
points de suspension, qui sont de lui.}

\page{40}
\emph{Remarque.} \struck{Supposons} Soit $S$ \uncertain{$= \mathrm{Spec}\,R$}
un schéma local noeth.\ complet, $\hat{X}$ un schéma formel propre
\struck{\ill{}} plat sur $\hat{S}$, tel que
$H^0(X_0, \mathcal{O}_{X_0}) \simeq k$ (\ill{} \ill{}). Supposons
\struck{$\underline{\mathrm{Pic}}^0$ soit lisse, soit que} aussi
\struck{\ill{}} que $\underline{\mathrm{Pic}}_{X_0/k}$ (p.ex.\
$\underline{\mathrm{Pic}}^0_{X_0/k}$ lisse) \ill{} \uncertain{tels} que
les $\underline{\mathrm{Pic}}^0_{X_n/S_n}$ \uncertain{soient} plats
\struck{soit que}, et formons
\[
  \underline{\mathrm{NS}}_{X_n/S_n} = \underline{\mathrm{Pic}}_{X_n/S_n} /
  \underline{\mathrm{Pic}}^0_{X_n/S_n} ,
\]
et
\[
  P_n = \underline{\mathrm{NS}}^{+}_{X_n/S_n} \quad
  (\text{polarisations \uncertain{relatives}}) .
\]
\ill{} \ill{} schémas \ill{}. Les $P_n$ se \struck{\ill{}}
\uncertain{recollent} \ill{} \ill{}, et \ill{} un schéma formel $\hat{P}$
sur $\hat{S}$, \ill{} \uncertain{définissant} un schéma \ill{} tel que
\[
  \hat{P} = \coprod_{\alpha} \hat{P}_{\alpha} ,
\]
les $P_{\alpha}$ étant finis sur $S$, et connexes. Procédant comme plus
haut, on trouve que les conditions suivantes sont
\uncertain{équivalentes} :

(i) $\exists$ $S' \to S$ \struck{\ill{}} \uncertain{propre} surjectif,
\struck{\ill{}} avec $\hat{X}_{\hat{S}'}$ algébrisable et \emph{projectif}
sur $S'$.

(ii) $\forall$ sous-schéma fermé $T$ de $S$, de dim $\leqslant 1$, si $T'$
est sa normalisée, $\hat{X}_{\hat{T}'}$ est algébrisable et projectif sur
$\hat{T}'$.

(iii) $\exists$ morphisme \uncertain{fini} surjectif $S' \to S$, tel que
$\hat{X}_{\hat{S}'}$ soit algébrisable et projectif.

(iv) Il existe un morphisme fini et surjectif $S' \to S$, tel que pour
toute composante irréductible $S'_i$ de $S'$, avec str.\ réduite induite,
$\hat{X}_{\hat{S}'_i}$ soit algébriquement projectif.
\struck{\ill{} \ill{} \ill{}} De \uncertain{plus}, \struck{\ill{} \ill{}
\ill{}} (iv), $\hat{X}_{\hat{S}}$ est algébrique.
\note{« $T \subset S$ » est ajouté au-dessus de « sous-schéma fermé » dans
(ii) ; « fini » au-dessus de « surjectif » dans (iii).}

\struck{(\ill{}) (iv) $\exists$ morphisme \ill{} fini \ill{} surjectif}
\note{la dernière ligne et la moitié de l'avant-dernière sont biffées ; un
crochet dans la marge les accompagne.}

\marginal{\textbf{N.B.} Cette condition \ill{} \ill{} telle \ill{} que si
\ill{} $\underline{\mathrm{Pic}}_{X_0/k}$ \ill{} \uncertain{lisse}) \ill{}
\ill{} \ill{} \ill{} $S' \to S$ \uncertain{fini et plat}. Car \ill{}
\ill{} \ill{} \uncertain{faisceaux} \ill{} \ill{}}
\note{note écrite de biais dans la marge gauche, reliée par une accolade aux
conditions (iii) et (iv).}

\marginal{\ill{} \uncertain{donc d'ailleurs} \uncertain{effectivité}}
\note{en bas de la marge gauche, en face de la ligne biffée.}

\end{document}
